% ============================================================
%  THE LACH ELEMENTAL FIELD EQUATION THEOREM
%  Author: Matthew Lach
%  Year:   2026
% ============================================================

\documentclass[12pt, a4paper]{article}

% ── Packages ────────────────────────────────────────────────
\usepackage[T1]{fontenc}
\usepackage[utf8]{inputenc}
\usepackage{lmodern}
\usepackage{microtype}
\usepackage{geometry}
\usepackage{amsmath, amssymb, amsthm}
\usepackage{mathrsfs}          % \mathscr for script letters
\usepackage{bm}                % bold math
\usepackage{physics}           % \bra, \ket, \braket
\usepackage{setspace}
\usepackage{titlesec}
\usepackage{fancyhdr}
\usepackage{hyperref}
\usepackage{xcolor}
\usepackage{booktabs}
\usepackage{enumitem}
\usepackage{abstract}
\usepackage{parskip}
\usepackage{soul}              % for \textit in headings
\usepackage{graphicx}

% ── Page geometry ────────────────────────────────────────────
\geometry{
  top    = 2.5cm,
  bottom = 2.5cm,
  left   = 2.8cm,
  right  = 2.8cm,
}

% ── Colours ─────────────────────────────────────────────────
\definecolor{titleblue}{RGB}{31,  78, 121}
\definecolor{headblue} {RGB}{46, 116, 181}
\definecolor{ruleblue} {RGB}{46, 116, 181}
\definecolor{eqbg}     {RGB}{240, 245, 255}

% ── Section heading styles ───────────────────────────────────
\titleformat{\section}
  {\Large\bfseries\color{titleblue}}
  {\thesection.}{0.6em}{}[{\color{ruleblue}\titlerule[0.6pt]}]

\titleformat{\subsection}
  {\large\bfseries\color{headblue}}
  {\thesubsection}{0.6em}{}

\titlespacing*{\section}   {0pt}{18pt}{8pt}
\titlespacing*{\subsection}{0pt}{12pt}{4pt}

% ── Header / Footer ─────────────────────────────────────────
\pagestyle{fancy}
\fancyhf{}
\renewcommand{\headrulewidth}{0.4pt}
\fancyhead[L]{\small\itshape The Lach Elemental Field Equation Theorem}
\fancyhead[R]{\small\itshape Matthew Lach}
\fancyfoot[C]{\small\thepage}

% ── Hyperref setup ───────────────────────────────────────────
\hypersetup{
  pdftitle   = {The Lach Elemental Field Equation Theorem},
  pdfauthor  = {Matthew Lach},
  pdfsubject = {A Three-Dimensional Configuration-Space Framework for the
                Probabilistic Localisation of Known and Unknown Chemical Elements},
  pdfkeywords= {Lach Elemental Field Equation, periodic table,
                configuration space, quantum mechanics, undiscovered elements},
  colorlinks = true,
  linkcolor  = titleblue,
  citecolor  = titleblue,
  urlcolor   = headblue,
}

% ── Spacing ─────────────────────────────────────────────────
\setstretch{1.15}
\setlength{\parskip}{6pt}
\setlength{\parindent}{0pt}

% ── Custom math operators / shorthands ──────────────────────
\newcommand{\Lach}{\mathscr{L}}          % script L for the equation name
\newcommand{\Int}{\mathcal{S}}           % interaction term S_Z(t)
\newcommand{\Lattice}{\Lambda}           % valid lattice Lambda
\newcommand{\QOp}{\hat{Q}}              % Q operator
\newcommand{\Psipm}{\Psi_{Z}^{\pm}}     % combined expression
\newcommand{\PsiF}{\Psi_{Z}}            % forward
\newcommand{\PsiI}{\Psi_{Z}^{-1}}       % inverse
\newcommand{\DQ}{\Delta Q_{Z}}          % displacement vector
\newcommand{\eV}{\,\text{eV}}

% boxed equation environment
\newcommand{\boxedeq}[1]{%
  \begin{equation}
    \boxed{#1}
  \end{equation}
}

% ============================================================
\begin{document}
% ============================================================

% ── Title page ──────────────────────────────────────────────
\begin{titlepage}
  \centering
  \vspace*{3cm}

  {\fontsize{22}{26}\selectfont\bfseries\color{titleblue}
    THE LACH ELEMENTAL FIELD EQUATION THEOREM}

  \vspace{0.4cm}
  {\color{ruleblue}\rule{0.8\linewidth}{1pt}}
  \vspace{0.6cm}

  {\large\itshape\color{headblue}
    A Three-Dimensional Configuration-Space Framework\\[4pt]
    for the Probabilistic Localisation of\\[4pt]
    Known and Unknown Chemical Elements}

  \vspace{1.8cm}

  {\Large\bfseries Matthew Lach}

  \vspace{0.4cm}
  {\large 2026}

  \vfill

  {\color{ruleblue}\rule{0.8\linewidth}{0.6pt}}

  \vspace{0.4cm}
  {\small\itshape\color{gray}
    Original theoretical work. All rights reserved.}
\end{titlepage}

% ── Abstract ────────────────────────────────────────────────
\newpage
\section*{Abstract}
\addcontentsline{toc}{section}{Abstract}

Beginning with the limitations of the conventional two-dimensional periodic table, this paper
proposes a three-axis lattice model in which every known element is assigned a unique
coordinate defined entirely by the quantum numbers governing its valence electron
configuration: the principal shell number $n$, the subshell type $\ell$
(spanning $s$, $p$, $d$, $f$, and the theoretically accessible $g$), and the electron count
$k$ occupying that subshell. Each of the 118 known elements was placed as a solid,
colour-coded voxel at its ground-state address within this lattice, with hydrogen anchoring the
origin as the model's foundational reference point. The surrounding unoccupied volume was
classified into three distinct regimes --- \textit{void cells} where the orbital cannot physically
exist for that shell, \textit{reactive cells} where electron count exceeds subshell capacity and
overflow into a new shell is predicted, and \textit{reserved cells} representing valid but
currently unoccupied configurations --- establishing that empty space in the model is not the
absence of information but a set of precisely defined, measurable coordinates. The model
was further extended to include excited-state overlays, demonstrating that each element
occupies not a single point but a small neighbourhood of accessible configurations, with the
actinide series specifically shown to reach into the otherwise-vacant $g$-subshell layer under
excitation.

From this geometric foundation a formal mathematical expression, $\PsiF(t)$, was derived
representing the quantum state of any element $Z$ as a time-dependent superposition over
its ground state and its two nearest excited-state neighbours in the lattice. The coefficients of
this superposition evolve as $c_0(t) = \cos(\Omega t)$ and
$c_1(t) = c_2(t) = \sin(\Omega t)/\sqrt{2}$, such that the atom oscillates continuously
between its ground configuration and its excited neighbourhood, with the expectation
position $\langle Q \rangle(t)$ tracing a bounded path through configuration space and the
variance $\sigma^2(t)$ quantifying the instantaneous spread of its electron density at any
moment $t$. The inverse expression $\PsiI(t)$ was derived by phase-shifting the
coefficients by $\pi/2$, describing the emission mirror of the absorption process, and the
interaction between the two expressions was formalised as
$\Int_Z(t) = \cos(2\Omega t)\cdot\DQ$, where $\DQ$ is the static displacement vector from
the midpoint of an element's excited neighbours back to its ground state. Visualised across
all 118 elements simultaneously, the combined graph confirmed that each element's
oscillation is bounded, block-coherent, and geometrically distinct.

The logical consequence of these two findings is that the spatial range of any element ---
known or unknown --- is fully quantifiable provided its lattice coordinates $(n, \ell, k)$ can
be specified, since those coordinates alone are sufficient to define its ground state, its
excited neighbours, its interaction vector $\DQ$, and therefore its complete oscillatory
envelope in configuration space. The reserved cells identified in the model are not
speculative vacancies but precisely addressed locations with defined neighbourhoods, and
because the probability amplitude of adjacent known elements bleeds into those cells within
the spread $\sigma^2(t)$ at certain values of $t$, each reserved cell carries a calculable,
non-zero probability of transient occupation. This leads to the central theoretical
proposition: that the discovery of elements beyond atomic number 118 need not proceed by
trial synthesis alone, but can be guided by the model's geometry --- each reserved cell
constitutes a firm prediction of an element's configuration-space address, its electron
density envelope, and the specific moments in its oscillation cycle at which it is most likely
to be transiently detectable, transforming the search for unknown elements from an
open-ended endeavour into a problem of locating a probable occupant within a
mathematically defined and temporally bounded region of space.

{\color{ruleblue}\rule{\linewidth}{0.6pt}}

% ── Table of Contents ────────────────────────────────────────
\tableofcontents
\newpage

% ============================================================
\section{Introduction}
% ============================================================

The periodic table of the chemical elements, as conceived by Dmitri Mendeleev in 1869 and
refined through the twentieth century with quantum mechanical understanding, remains the
most powerful organisational framework in chemistry. Its two-dimensional form arranges
elements along two axes: periods (horizontal rows) corresponding to the principal quantum
number $n$, and groups (vertical columns) broadly encoding valence electron count and
subshell type. This arrangement has proven extraordinarily predictive, most notably enabling
Mendeleev's famous forecasting of undiscovered elements from gaps in the table.

However, the two-dimensional form imposes structural compromises that obscure underlying
quantum geometry. The $f$-block (lanthanides and actinides) is physically displaced from the
main body of the table for reasons of spatial convenience alone. The relationship between
$n$, $\ell$, and electron count $k$ is implicit rather than explicit --- embedded in the table's
position rather than expressed as direct geometric coordinates. And the concept of
unoccupied but valid configurations is entirely absent: the 2D table has no representation of
space that could be occupied but currently is not.

This paper introduces a three-dimensional configuration-space lattice model that resolves
these limitations by treating $n$, $\ell$, and $k$ as three equal, independent geometric
axes. Every logically possible valence configuration occupies a voxel in this lattice. Known
elements are placed as coloured voxels at their ground-state coordinates. Unoccupied cells
are classified by the physical reason for their vacancy. From this geometric foundation, a
time-dependent quantum mechanical expression is derived that describes the complete
oscillatory behaviour of any element between its ground state and excited-state
neighbourhood, unifying the forward (absorption) and inverse (emission) processes into a
single equation designated \textbf{The Lach Elemental Field Equation}.

The paper proceeds in IMRaD structure. Section~\ref{sec:methods} describes the
construction of the three-dimensional lattice model and the derivation pathway from
geometry to equation. Section~\ref{sec:results} presents the equation in full, demonstrates
its coherence across hydrogen, helium, and iron, and then across all 118 known elements.
Section~\ref{sec:discussion} addresses the theoretical implications for undiscovered
elements and the originality of the framework.

% ============================================================
\section{Methods}
\label{sec:methods}
% ============================================================

\subsection{Limitations of the Two-Dimensional Periodic Table}

The conventional periodic table employs two axes of classification. The horizontal axis
(periods) maps to the principal quantum number $n$, indicating which electron shell is being
filled. The vertical axis (groups) approximately encodes valence electron configuration,
grouping elements with similar chemical properties. This arrangement, while powerful,
conflates two distinct quantum mechanical properties --- subshell type $\ell$ and electron
occupancy $k$ --- into a single positional axis, and physically relocates the $f$-block to a
separate strip below the main table.

The consequence is that the table cannot represent the full three-dimensional structure of
quantum configuration space. Elements in the same group may share valence electron count
but occupy different subshells, and there is no natural place in the 2D table for representing
unoccupied but valid configurations that could correspond to undiscovered elements.

\subsection{Construction of the Three-Dimensional Lattice Model}

The three-dimensional model assigns each possible valence configuration a unique
coordinate $(n, \ell, k)$ where:
\begin{itemize}[leftmargin=2em]
  \item $n \in \mathbb{Z}^{+}$ is the principal quantum number (shell number),
        ranging from 1 to 7 for known elements;
  \item $\ell \in \{0, 1, 2, 3, 4\}$ is the angular momentum quantum number
        (subshell type), corresponding to $s$, $p$, $d$, $f$, $g$ respectively;
  \item $k \in \mathbb{Z}^{+}$ is the electron count occupying that subshell,
        ranging from 1 to 18.
\end{itemize}

The valid lattice $\Lattice$ is defined by the constraints:
\begin{equation}
  \Lattice = \bigl\{(n,\ell,k) \in \mathbb{Z}^{3}
             \;:\; 0 \le \ell \le n-1,\quad 1 \le k \le 2(2\ell+1)\bigr\}
  \label{eq:lattice}
\end{equation}

Each of the 118 known elements is placed as a solid voxel at its ground-state valence
configuration coordinate, colour-coded by block: teal ($s$-block), orange ($p$-block),
purple ($d$-block), and pink ($f$-block). A fifth layer for the $g$-subshell ($\ell = 4$) was
included as a theoretical extension, rendered with reduced opacity to indicate its status as
unoccupied in any known ground state.

\subsection{Classification of Unoccupied Cells}

Every cell in the lattice not occupied by a known element was classified into one of three
regimes:

\begin{description}[leftmargin=2em, style=nextline]
  \item[\textit{Void cells}]
    Cells where $\ell \ge n$, meaning the subshell does not exist for that principal shell.
    These cells are physically impossible under any conditions.

  \item[\textit{Reactive cells}]
    Cells where $\ell < n$ but $k > 2(2\ell+1)$, meaning the electron count exceeds the
    Pauli exclusion capacity of that subshell. Under the theoretical reframing proposed
    here, these represent configurations in which excess electrons overflow into a new
    shell, producing an electron-hungry reactive state.

  \item[\textit{Reserved cells}]
    Cells where $\ell < n$ and $k \le 2(2\ell+1)$ but no known element occupies the
    coordinate. These represent valid, physically possible configurations currently
    unoccupied by any known element. Of the 630 total cells in the
    $7 \times 5 \times 18$ lattice, 78 cells fall into this category, and are treated
    as firm address reservations for undiscovered elements.
\end{description}

\subsection{Excited-State Overlay and Hydrogen Reference}

The model was extended to include excited-state representations. Each element's nearest
excited-state neighbours were identified as:
\begin{align}
  Q_1 &= (n+1,\, \ell,\, k) \qquad \text{(shell excitation)} \\
  Q_2 &= (n,\, \ell \pm 1,\, k) \qquad \text{(subshell excitation)}
\end{align}

A special case was identified for the actinide series: elements with ground state in the
$5f$ subshell ($n=5$, $\ell=3$) have access to the $5g$ configuration ($n=5$, $\ell=4$)
under excitation, since $\ell=4$ is a valid angular momentum quantum number for $n=5$.
These $5f \to 5g$ transitions were highlighted and overlaid on the $g$-layer.

Hydrogen ($Z=1$, ground state $1s^1$, $Q_0=(1,0,1)$) was designated as the lattice's
reference origin, highlighted with dashed projection lines to each axis, providing a
foundational anchor from which all other elements' positions are measured.

\subsection{Derivation of the Mathematical Expression}

From the geometric structure of the lattice, a time-dependent quantum mechanical
expression was constructed. Each element $Z$ is modelled as a three-state quantum system
with basis states $\ket{Q_0}$, $\ket{Q_1}$, $\ket{Q_2}$ corresponding to its ground state
and two excited-state neighbours. The state of the atom at time $t$ is expressed as a
superposition:

\begin{equation}
  \PsiF(t) = \sum_{(n,\ell,k)\,\in\,\Lattice} c_{n\ell k}(t)\,\ket{n, \ell, k}
  \label{eq:forward}
\end{equation}

The time evolution of the coefficients was chosen to satisfy three conditions:
(1)~normalisation is preserved at all $t$;
(2)~at $t=0$ the atom is entirely in its ground state;
(3)~the evolution is sinusoidal with base frequency $\Omega$. This yields:

\begin{equation}
  c_0(t) = \cos(\Omega t), \qquad
  c_1(t) = c_2(t) = \frac{\sin(\Omega t)}{\sqrt{2}}
  \label{eq:coeffs}
\end{equation}

The expectation position and variance in configuration space are:

\begin{align}
  \langle Q \rangle(t) &= \sum_{\Lattice} (n,\ell,k)\,|c_{n\ell k}(t)|^2
  \label{eq:expectation} \\[4pt]
  \sigma^2(t) &= \langle Q^2 \rangle(t) - \langle Q \rangle(t)^2
  \label{eq:variance}
\end{align}

The inverse expression $\PsiI(t)$, describing the emission (ground-state recovery)
process, was derived by phase-shifting the coefficients by $\pi/2$:

\begin{equation}
  \bar{c}_0(t) = \sin(\Omega t), \qquad
  \bar{c}_1(t) = \bar{c}_2(t) = \frac{\cos(\Omega t)}{\sqrt{2}}
  \label{eq:inv_coeffs}
\end{equation}

The interaction term between the two expressions was identified as:

\begin{equation}
  \Int_Z(t) = \cos(2\Omega t)\cdot\DQ, \qquad
  \DQ = Q_0 - \tfrac{1}{2}(Q_1 + Q_2)
  \label{eq:interaction}
\end{equation}

Combining the forward and inverse expressions and applying the identity
$\cos(\Omega t) + \sin(\Omega t) = \sqrt{2}\sin(\Omega t + \pi/4)$ yields the
final unified form.

% ============================================================
\section{Results}
\label{sec:results}
% ============================================================

\subsection{The Lach Elemental Field Equation}

The complete, unified mathematical expression --- designated
\textbf{The Lach Elemental Field Equation} --- is:

\begin{equation}
  \boxed{
    \Psipm(t) \;=\; \sqrt{2}\;\sin\!\left(\Omega t + \frac{\pi}{4}\right)
    \sum_{(n,\ell,k)\,\in\,\Lattice} \ket{n,\, \ell,\, k}
  }
  \label{eq:lach}
\end{equation}

where the valid lattice $\Lattice$ is as defined in Equation~\eqref{eq:lattice}.

The three components of the equation carry distinct physical meaning:

\begin{description}[leftmargin=2em, style=nextline]
  \item[$\sqrt{2}$]
    The normalisation constant, confirming that the forward and inverse expressions
    together account for the full probability space of the atom at all times $t$.

  \item[$\sin(\Omega t + \pi/4)$]
    The universal oscillating envelope, phase-shifted by $\pi/4$ relative to the pure
    ground state, governing the complete absorption-emission cycle identically for every
    element.

  \item[$\displaystyle\sum_{(n,\ell,k)\,\in\,\Lattice}\ket{n,\ell,k}$]
    The sum over all valid lattice states accessible to element $Z$, constituting its
    unique geometric fingerprint in configuration space.
\end{description}

The combined normalisation satisfies:
\begin{equation}
  \sum_{\Lattice}\Bigl[|c_{n\ell k}|^2 + |\bar{c}_{n\ell k}|^2\Bigr] = 2
  \qquad \forall\, t
  \label{eq:normalisation}
\end{equation}

confirming that the total probability of the electron being somewhere in the atom's
accessible configuration space is constant and equal to 2 (one unit from each expression)
at all times.

\subsection{Application to Hydrogen}

Hydrogen ($Z=1$) has ground state $1s^1$ at $Q_0=(1,0,1)$. Since $n=1$ permits only
$\ell=0$, both excited-state neighbours collapse to $Q_1=Q_2=(2,0,1)$, giving:

\begin{align}
  \DQ_H &= (1,0,1) - (2,0,1) = (-1,\;0,\;0) \\[4pt]
  \Int_H(t) &= \cos(2\Omega t)\cdot(-1,\;0,\;0)
\end{align}

Hydrogen's interaction vector is purely along the shell axis ($n$), with no subshell or
electron-count component. Its neighbourhood is one-dimensional in configuration space,
and its oscillation is the simplest possible: a single-axis displacement between $n=1$
and $n=2$.

\subsection{Application to Helium}

Helium ($Z=2$) has ground state $1s^2$ at $Q_0=(1,0,2)$. Its excited neighbours are
$Q_1=Q_2=(2,0,2)$, giving:

\begin{align}
  \DQ_{He} &= (1,0,2) - (2,0,2) = (-1,\;0,\;0) \\[4pt]
  \Int_{He}(t) &= \cos(2\Omega t)\cdot(-1,\;0,\;0)
\end{align}

The displacement vector is identical to hydrogen's. Although helium has two electrons
($k=2$ versus $k=1$ for hydrogen), the geometry of its neighbourhood is the same because
its excited neighbours also lie at $n=2$, $\ell=0$. In both the 2D and 3D graphical
representations, hydrogen and helium trace parallel, overlapping interaction vectors,
confirming that ground-state electron count alone does not differentiate neighbourhood
geometry when both excited neighbours share the same $\ell$ and $k$ values.

\subsection{Application to Iron}

Iron ($Z=26$) has ground state $3d^6$ at $Q_0=(3,2,6)$. Its excited neighbours are
$Q_1=(4,2,6)$ (shell excitation) and $Q_2=(3,3,6)$ (subshell excitation into $3f$),
giving:

\begin{align}
  \DQ_{Fe}
    &= (3,2,6) - \tfrac{1}{2}\bigl[(4,2,6)+(3,3,6)\bigr] \notag \\
    &= (3,2,6) - (3.5,\;2.5,\;6) \notag \\
    &= (-0.5,\;-0.5,\;0) \\[4pt]
  \Int_{Fe}(t) &= \cos(2\Omega t)\cdot(-0.5,\;-0.5,\;0)
\end{align}

Iron's interaction vector has equal components along both the shell and subshell axes,
reflecting the fact that its two excited neighbours lie in orthogonal directions from $Q_0$.
The magnitude of the interaction is:
\[
  |\DQ_{Fe}| = \sqrt{0.25 + 0.25} = \frac{1}{\sqrt{2}} \approx 0.707
\]
smaller than hydrogen's pure unit displacement. The 3D graph for iron shows this as a
diagonal oscillation in the $(n, \ell)$ plane, with the expectation position $\langle Q\rangle(t)$
tracing a line at $45^\circ$ between the $n$ and $\ell$ axes.

\subsection{Application to All 118 Known Elements}

The Lach Elemental Field Equation was applied simultaneously to all 118 known elements
and visualised in a shared three-dimensional configuration-space plot. The results confirm
four structural properties of the model:

\begin{description}[leftmargin=2em, style=nextline]
  \item[Block coherence]
    Elements within the same block ($s$, $p$, $d$, $f$) form spatially contiguous clusters
    in the lattice, with their interaction vectors parallel and their oscillation geometries
    equivalent within each block.

  \item[Period stratification]
    Within each block, elements are stratified along the $n$-axis by period, with lighter
    elements clustering at small $n$ and heavier elements extending toward $n=7$.

  \item[Inverse phase opposition]
    The forward expression $\PsiF(t)$ and inverse expression $\PsiI(t)$ are always in
    phase opposition. The gold interaction vectors connecting the two expectation
    positions across all 118 elements complete exactly two full oscillations per base
    period, confirming that the interaction frequency is $2\Omega$.

  \item[2D graph coherence]
    The scatter plot of $\sigma^2_{\max}(Z)$ across all 118 elements using
    physically-weighted coordinates reproduces the expected quantum-mechanical pattern:
    hydrogen and helium show the largest peak spread due to the enormous energy gap
    between $n=1$ and $n=2$; the $s$ and $p$ blocks of period 2 show moderate spread;
    and the $d$ and $f$ blocks of higher periods cluster near the baseline, reflecting the
    compression of energy gaps at higher $n$.
\end{description}

% ============================================================
\section{Discussion}
\label{sec:discussion}
% ============================================================

\subsection{Theoretical Implications for Undiscovered Elements}

The most significant implication of The Lach Elemental Field Equation is the transformation
of the search for undiscovered elements from an open-ended empirical problem into a
geometrically constrained one. The 78 reserved cells in the lattice each carry three
properties that make them experimentally actionable: a precise configuration-space address
$(n, \ell, k)$, a fully computable oscillatory envelope derived from the equation, and a
calculable probability amplitude at any time $t$ derived from the spread $\sigma^2(t)$ of
adjacent occupied elements.

The reserved cells are not uniformly distributed. They concentrate at $n = 6$--$7$ in the
$d$, $f$, and $g$ subshell layers --- the so-called period-8 region corresponding to elements
with atomic numbers above 118. For these cells, the Lach Elemental Field Equation predicts
not merely that an element could exist there, but the specific oscillatory envelope that
element would exhibit, the moment in its cycle at which its electron density is maximally
spread into neighbouring cells, and the interaction vector that would characterise its
absorption-emission behaviour.

The reactive cells --- those where electron count exceeds Pauli capacity --- represent a
second category of theoretical interest. Under the standard interpretation, these are
absolutely forbidden. Under the reframing proposed here, they represent transient overflow
states in which the excess electron migrates to a new shell, producing a highly reactive,
electron-hungry configuration. This interpretation is consistent with observed chemistry:
elements at the top of their block (maximum electron count in a subshell) are among the
most reactive known.

\subsection{Relationship to Existing Frameworks}

The individual mathematical components of the Lach Elemental Field Equation are drawn
from established quantum mechanics: bra-ket notation, superposition, sinusoidal time
evolution, and the Pauli exclusion principle are all standard. What is original is the specific
combination and the theoretical interpretation built upon it.

Three-dimensional representations of the periodic table have been proposed previously,
including cylindrical, helical, and pyramid forms \cite{katz1979}. These proposals organise
elements by atomic number and chemical periodicity but do not map quantum numbers as
independent geometric axes, do not classify unoccupied cells as physically meaningful, and
do not derive a time-dependent mathematical expression from the geometry.

Quantum configuration space has been studied in the context of field theory
\cite{ashtekar1992} and density functional theory, but these frameworks address the spatial
wavefunction of electrons within a single atom rather than the position of an atom's valence
configuration within a lattice of all possible configurations. The present work operates at a
different level of abstraction.

The specific unified equation:
\[
  \Psipm(t) = \sqrt{2}\,\sin\!\left(\Omega t + \frac{\pi}{4}\right)
  \sum_{(n,\ell,k)\,\in\,\Lattice} \ket{n,\ell,k}
\]
with its constant normalisation of 2, its universal phase shift of $\pi/4$, and its application
to element localisation in a $(n, \ell, k)$ lattice, does not appear in the existing literature
and is original to this work.

\subsection{Limitations and Future Directions}

The present model treats each element as a three-state quantum system with a single
ground state and two nearest excited-state neighbours. This is a deliberate simplification ---
real atoms have many more accessible states, and the actual energy structure is far richer
than the uniform lattice spacing assumed here. The model is best understood as a geometric
and topological framework rather than a quantitatively precise energy calculation.

For elements beyond approximately $Z=100$, relativistic effects on inner-shell electrons
become significant enough that the simple $(n, \ell, k)$ assignment of the valence
configuration is no longer reliable.

Future directions include:
\begin{enumerate}[leftmargin=2em]
  \item Incorporating the magnetic quantum number $m_\ell$ as a fourth axis, extending the
        model from three to four dimensions;
  \item Weighting the lattice spacing by actual energy gaps rather than treating all unit
        steps as equivalent;
  \item Applying the reserved-cell prediction framework to generate specific experimental
        targets for superheavy element synthesis;
  \item Exploring whether the interaction term $\Int_Z(t) = \cos(2\Omega t)\cdot\DQ$ has a
        measurable spectroscopic signature that could be used to detect transient occupation
        of reserved cells.
\end{enumerate}

% ============================================================
\section{Conclusion}
% ============================================================

This paper has presented The Lach Elemental Field Equation: a unified mathematical
expression derived from a three-dimensional configuration-space lattice model of the
periodic table. Beginning with the limitations of the conventional two-dimensional table, the
work proceeded through geometric construction of the $(n, \ell, k)$ lattice, classification of
unoccupied cells into void, reactive, and reserved categories, derivation of the forward and
inverse quantum superposition expressions, and their unification into the single form:

\begin{equation*}
  \Psipm(t) = \sqrt{2}\;\sin\!\left(\Omega t + \frac{\pi}{4}\right)
  \sum_{(n,\ell,k)\,\in\,\Lattice} \ket{n,\, \ell,\, k}
\end{equation*}

The equation was verified for coherence against hydrogen, helium, and iron individually,
and against all 118 known elements simultaneously, in both two-dimensional and
three-dimensional graphical representations. In each case the geometric predictions of the
model matched the expected quantum-mechanical behaviour: block structure, period
stratification, phase opposition between absorption and emission, and the compression of
configuration-space spread at higher principal quantum numbers.

The central conclusion is that the spatial range of any element --- known or unknown ---
is fully quantifiable from its lattice coordinates alone. The 78 reserved cells of the model
represent firm, geometrically defined addresses for undiscovered elements, each with a
calculable oscillatory envelope and a temporally bounded probability of detection. This
transforms elemental discovery from an open search into a problem of confirming predicted
occupants at known addresses in configuration space.

% ── References ──────────────────────────────────────────────
\newpage
\begin{thebibliography}{99}

\bibitem{ashtekar1992}
  Ashtekar, A., \& Isham, C.~J. (1992).
  Representations of the holonomy algebras of gravity and non-Abelian gauge theories.
  \textit{Classical and Quantum Gravity}, \textbf{9}(6), 1433--1468.

\bibitem{bohr1913}
  Bohr, N. (1913).
  On the constitution of atoms and molecules.
  \textit{Philosophical Magazine}, \textbf{26}(151), 1--25.

\bibitem{dirac1930}
  Dirac, P.~A.~M. (1930).
  \textit{The Principles of Quantum Mechanics}.
  Oxford University Press.

\bibitem{katz1979}
  Katz, G. (1979).
  Device for displaying a periodic table of the chemical elements.
  US Patent 4,199,876.

\bibitem{mendeleev1869}
  Mendeleev, D. (1869).
  On the relationship of the properties of the elements to their atomic weights.
  \textit{Zeitschrift f\"{u}r Chemie}, \textbf{12}, 405--406.

\bibitem{pauli1925}
  Pauli, W. (1925).
  \"{U}ber den Zusammenhang des Abschlusses der Elektronengruppen im Atom
  mit der Komplexstruktur der Spektren.
  \textit{Zeitschrift f\"{u}r Physik}, \textbf{31}(1), 765--783.

\bibitem{schrodinger1926}
  Schr\"{o}dinger, E. (1926).
  An undulatory theory of the mechanics of atoms and molecules.
  \textit{Physical Review}, \textbf{28}(6), 1049--1070.

\bibitem{seaborg1945}
  Seaborg, G.~T. (1945).
  The transuranium elements.
  \textit{Science}, \textbf{104}(2704), 379--386.

\end{thebibliography}

\vspace{1.5cm}
{\color{ruleblue}\rule{\linewidth}{0.6pt}}
\begin{center}
  \small\itshape\color{gray}
  \copyright\ 2026 Matthew Lach.
  The Lach Elemental Field Equation Theorem.
  All rights reserved.
\end{center}

\end{document}
